%=====================================================================
\chapter{The Stack, and Your Shop in One Command}\label{ch:stack}
%=====================================================================
\locator{2}{Part II --- Standing the Shop Up}

\begin{outcomes}
By the end of this chapter you will be able to:
\begin{tight}
  \item \textbf{Name} the four layers a web shop runs on and \textbf{state}
        what each one is responsible for.
  \item \textbf{Explain} what a container is, and \textbf{distinguish} it from
        a virtual machine and from an installer.
  \item \textbf{Read} a Compose file and \textbf{say} what each service, port
        and volume in it does.
  \item \textbf{Bring up} a working storefront, \textbf{verify} it from outside,
        and \textbf{reset} it to a known state.
  \item \textbf{Measure} what the stack costs to start and to store, and
        \textbf{distinguish} the cost paid once from the cost paid daily.
\end{tight}
\end{outcomes}

\begin{prereq}
\chref{planning}: the decision to use a self-hosted open-source platform, and
why. Appendix~A must be finished before this chapter --- you need a working
container runtime, and if it is going to fight you, it will fight you now
rather than in Week~9.
\end{prereq}

\begin{keyterms}
stack $\bullet$ LAMP $\bullet$ web server $\bullet$ application server
$\bullet$ database server $\bullet$ container $\bullet$ image $\bullet$ volume
$\bullet$ port mapping $\bullet$ Docker Compose $\bullet$ service $\bullet$
environment variable $\bullet$ reproducibility $\bullet$ idempotent
$\bullet$ pinning $\bullet$ checksum
\end{keyterms}

\section{Four Layers, and What Each One Does}

When you type a shop's address into a browser, four distinct pieces of software
answer in turn. Knowing which is which is the difference between fixing a
problem and restarting things until it goes away.

\dsfig{%
\begin{tikzpicture}[
  lyr/.style={rounded corners=3pt, draw=#1, line width=1.1pt, fill=#1!8,
              text=#1!60!black, font=\bfseries\small, align=center,
              text width=3.1cm, minimum height=1.15cm},
  job/.style={font=\scriptsize, text=neutral, align=left, text width=7.4cm,
              anchor=west},
  arr/.style={-{Stealth[length=2.2mm]}, primary!50, line width=0.9pt}
]
  \node[lyr=neutral]   (os) at (0,0)    {Operating system\\{\scriptsize
       \mdseries Debian Linux}};
  \node[lyr=secondary] (ws) at (0,1.55) {Web server\\{\scriptsize\mdseries
       Apache 2.4}};
  \node[lyr=greenhl]   (ap) at (0,3.10) {Application\\{\scriptsize\mdseries
       PHP 8.2 + OpenCart}};
  \node[lyr=purplehl]  (db) at (0,4.65) {Database\\{\scriptsize\mdseries
       MariaDB 11}};

  \foreach \a/\b in {os/ws, ws/ap, ap/db}{\draw[arr] (\a) -- (\b);}

  \node[job] at (2.0,0)    {Files, processes, the network socket. You will
       almost never touch it.};
  \node[job] at (2.0,1.55) {Accepts the HTTP request, decides which file
       answers it, serves images and stylesheets itself.};
  \node[job] at (2.0,3.10) {Runs the shop's code: works out what the page
       should say, asks the database, builds the HTML.};
  \node[job] at (2.0,4.65) {Remembers everything --- products, prices, baskets,
       orders, customers. The only layer whose loss is unrecoverable.};

  \node[font=\scriptsize\itshape, text=accent, align=center] at (4.8,5.85)
       {A slow page is usually this layer. A lost order is always this one.};
\end{tikzpicture}}%
{The four layers. The traditional name for this arrangement is LAMP --- Linux,
Apache, MySQL, PHP --- and although every letter has alternatives now, the
shape has not changed in twenty-five years.}{lampstack}

\begin{conceptbox}[Which layer is this problem in?]
A diagnostic habit worth building on day one, because it turns ``the site is
broken'' into a question with an answer.

\begin{tight}
  \item \textbf{Nothing answers at all}, the browser cannot connect: the web
        server is down, or the port is wrong. Nothing above it has run.
  \item \textbf{A page appears but says ``500''}: the web server is fine; the
        application failed. The reason is in the PHP error log, never on the
        screen --- and \chref{security} explains why you want it that way.
  \item \textbf{A page appears, looks right, but data is missing or stale}: the
        application ran and the database answered something unexpected.
  \item \textbf{Everything works and is slow}: usually the database, and
        \chref{catalogue} measures exactly that.
\end{tight}
\end{conceptbox}

\section{Containers, and What They Are Not}

The traditional way to get those four layers is to install them: a web server, a
database, a PHP runtime, each with its own installer, its own configuration
file in its own place, and its own way of being slightly different on your
machine from on everybody else's.

A \term{container} replaces that with a package that already contains the
software \emph{and} everything it depends on, down to the system libraries,
arranged exactly as its author arranged it.

\begin{definitionbox}[Image, Container, Volume]
An \term{image} is a read-only, layered filesystem containing an application
and its dependencies. It is built once and is identical everywhere.

\smallskip
A \term{container} is a running instance of an image. It has its own filesystem
view and its own network identity, but it shares the host's kernel --- which is
what makes it start in a second rather than a minute.

\smallskip
A \term{volume} is storage that lives \emph{outside} the container and survives
it. This is the one that matters: everything inside a container is thrown away
when the container is removed, so anything you want to keep --- your database,
your shop --- must be in a volume.
\end{definitionbox}

\begin{alertbox}[Three things a container is not]
\textbf{It is not a virtual machine.} A VM carries a whole operating system
including its own kernel, takes a minute to boot and gigabytes of memory. A
container shares the host kernel and starts in about a second. The isolation is
correspondingly weaker, which matters in production and not in this course.

\smallskip
\textbf{It is not a backup.} ``It is in the container'' means it will be gone
when the container is replaced, which happens every time you rebuild. Your shop
survives because it is in a \emph{volume}, and \chref{deploy} is about what
happens when somebody has not understood that distinction.

\smallskip
\textbf{It is not automatically secure.} A container with a weak password is a
weak password that starts quickly. \chref{security} attacks this very stack and
finds things.
\end{alertbox}

\section{The Three Services in This Course's Stack}

The whole stack is described by one file, \texttt{code/docker-compose.yml},
and it declares three \term{services}.

\begin{lstlisting}[language=yaml]
services:
  db:                              # MariaDB 11: remembers everything
    image: mariadb:11
    volumes:
      - dbdata:/var/lib/mysql      # the shop's memory, kept outside
    healthcheck: ...               # "accepts a login", not "started"

  shop:                            # Apache + PHP 8.2 + OpenCart 4.1.0.4
    build: ./shop
    depends_on:
      db: { condition: service_healthy }
    ports:
      - "8090:80"                  # your 8090 -> the container's 80
    volumes:
      - shopdata:/var/www/html

  adminer:                         # a database console in the browser
    image: adminer:latest
    ports:
      - "8091:8080"
\end{lstlisting}

Three decisions in that file are worth understanding, because each one is a
mistake people make.

\textbf{The healthcheck, not just \texttt{depends\_on}.} Ordering container
\emph{starts} is not enough. MariaDB accepts TCP connections several seconds
before it will accept a login, so a shop that waits only for the database to
have started will try to install itself against a database that refuses it ---
intermittently, and more often on a fast machine. The condition
\texttt{service\_healthy} waits for the database to say it is actually ready.

\textbf{The port mapping is a translation, not a setting.} \texttt{8090:80}
means ``traffic arriving at port 8090 on my machine goes to port 80 inside the
container''. The shop always believes it is on port 80. If 8090 is busy on your
machine --- and on a university laptop it often is --- change the left-hand
number only.

\textbf{Named volumes, so that a reset is deliberate.} \texttt{docker compose
down} stops everything and keeps your shop. \texttt{docker compose down -v}
removes the volumes too, and the next start reinstalls from nothing. That
second command is how you recover from a shop you have broken beyond
understanding, and it is also how you destroy a semester's work in week 14. It
is worth knowing which is which before you need either.

\begin{conceptbox}[Why the image is built here rather than downloaded]
Most teaching material says ``pull the OpenCart image''. There is no longer one
to pull: there has never been an official OpenCart image, and the widely-used
third-party one was withdrawn. Every such instruction on the web is now dead.

\smallskip
So \texttt{code/shop/Dockerfile} builds it, from the official PHP image plus a
\emph{pinned} OpenCart release --- pinned by version and by SHA-256 checksum,
so that the shop you build this year is the shop a classmate builds next year.
Two lines in that file do the real work:

\begin{tight}
  \item \texttt{ARG OC\_VERSION=4.1.0.4} --- the exact release.
  \item \texttt{echo "\$\{OC\_SHA256\}  /tmp/oc.zip | sha256sum -c -} --- the
        download really is the file we tested against. A version number alone
        is not reproducible, because a release asset can be replaced.
\end{tight}

\smallskip
There is a second reason, and it is the pedagogical one: \dsref{lampstack} is a
diagram of something, and pulling one opaque image would have left you unable
to point at any of it. Here the layers are in the file, in order.
\end{conceptbox}

\begin{alertbox}[OpenCart ships its own Dockerfile, and it does not work]
The release archive contains a \texttt{Dockerfile} and a
\texttt{docker-compose.yml} at the top level. This course uses neither, for two
reasons you can verify yourself in five minutes.

\smallskip
The Dockerfile contains a shell syntax error --- \texttt{docker-php-ext-install
zip \&\& \&\& docker-php-ext-enable zip} --- so it has never built. The compose
file pins \texttt{mysql:5.7}, which reached end of life in October~2023 and
receives no further security fixes, and starts Redis, Memcached and PostgreSQL
that a stock install does not use.

\smallskip
The lesson is not that OpenCart is bad software; it is a mature project that
runs a large share of the world's small shops. The lesson is that an artefact
being official is evidence and not authority. You are allowed to check, and in
this course you are expected to.
\end{alertbox}

\section{What Happens the First Time}

One command brings the whole thing up:

\begin{lstlisting}[language=bash]
docker compose up -d --build
\end{lstlisting}

On a first run this does rather a lot: downloads about 1.5~GB of images, builds
the shop image, starts three containers, waits for the database, and then runs
OpenCart's command-line installer, which creates 159 tables and populates them
with a demonstration catalogue. On every later run it does almost none of it.

That difference is worth measuring rather than asserting.

\begin{measured}[What the Stack Costs to Start]
Produced by \texttt{code/m05\_coldstart.py --scratch}, which removes every
image and volume first, on the machine in Appendix~A --- 4 cores, 7.7~GiB
available to the container runtime.

\rowcolors{2}{white}{lightbg}
\begin{longtable}{@{}L{5.8cm}L{2.6cm}L{6.0cm}@{}}
\toprule
\rowcolor{primary!15}
\textbf{Start} & \textbf{To first 200} & \textbf{What it had to do} \\
\midrule
\endhead
From scratch (no images at all) & \textbf{338.4 s} & Download 1.5 GB, compile
three PHP extensions, fetch and verify OpenCart, install 159 tables. Of this,
214.6 s was the build alone \\
Cold (images built, no shop) & \textbf{34.1 s} & Start three containers, wait
for the database, install OpenCart \\
Warm (the shop already exists) & \textbf{16.4 s} & Start three containers.
Nothing else \\
\bottomrule
\end{longtable}

\smallskip
\textbf{Footprint.} Images: shop 882~MB, MariaDB 462~MB, Adminer 164~MB ---
\textbf{1{,}508~MB} in all. Volumes after installation: database 181~MB, shop
files 40~MB. About 1.7~GB of disk for a complete shop.

\smallskip
\textbf{The ratio is the point.} From scratch to warm is \textbf{20.6$\times$}.
You pay the 338 seconds once, in Week~3, and the 16 seconds perhaps sixty times
over the semester. Quoting a single ``startup time'' for this stack would
describe an experience nobody has.

\smallskip
\textbf{On variance.} An earlier run of the same measurement on the same
machine gave 17.6 s cold and 7.7 s warm --- roughly half --- because less was
running at the time. Your numbers will differ from these by more than you
expect. The \emph{ratio} between cold and warm should not.
\end{measured}

\section{Looking At What You Were Given}

The installer leaves you a working shop with a demonstration catalogue. It is
worth knowing its shape, because every chapter in Part~II changes part of it.

\begin{examplebox}[The shop, counted]
Asked of the database directly, through Adminer or the command line:

\begin{lstlisting}[language=sql]
SELECT (SELECT COUNT(*) FROM oc_product)    AS products,
       (SELECT COUNT(*) FROM oc_category)   AS categories,
       (SELECT COUNT(*) FROM oc_order)      AS orders,
       (SELECT COUNT(*) FROM oc_customer)   AS customers;
\end{lstlisting}

\begin{tight}
  \item \textbf{159 tables} in total.
  \item \textbf{19 products} and \textbf{38 categories} of demonstration data.
  \item \textbf{0 orders} and \textbf{0 customers} --- nobody has shopped here
        yet.
  \item \textbf{18 tables whose name begins \texttt{oc\_product}}. One product
        is not one row, and \chref{catalogue} is about why.
\end{tight}

\smallskip
That last number is the one to sit with. A beginner expects a products table.
There are eighteen, because a product has descriptions in several languages,
images, options, attributes, discounts, related items and category memberships
--- and each of those is a different relationship.
\end{examplebox}

\begin{pitfall}
\begin{tight}
  \item \textbf{Running \texttt{down -v} when you meant \texttt{down}.} The
        \texttt{-v} removes the volumes, which is your shop. There is no undo.
  \item \textbf{Editing files inside a running container.} They are gone at the
        next rebuild. Changes that must survive belong in the image or in a
        volume.
  \item \textbf{Assuming \texttt{up} means ready.} It means \emph{started}.
        On a first run the shop answers about twenty seconds later; every
        measurement in this book calls \texttt{waitfor.py} first for exactly
        this reason.
  \item \textbf{Changing the right-hand side of a port mapping.} In
        \texttt{8090:80}, the 80 is the port inside the container and the shop
        expects it. Change the left number.
  \item \textbf{Treating a container as a backup.} It is the opposite of one.
  \item \textbf{Giving up when it does not start.} Appendix~A lists the five
        failures that account for nearly all of them, and four are a busy port
        or a runtime that is not actually running.
\end{tight}
\end{pitfall}

\begin{lab}[Your shop, up, verified, broken and restored]
From here to \chref{project} this is the same shop. Treat it accordingly.

\begin{lstlisting}[language=bash]
# --- 1. Bring the whole stack up -----------------------------------
# First run: expect about five minutes and 1.5 GB of download.
cd code
cp .env.example .env
docker compose up -d --build

# --- 2. Watch it install -------------------------------------------
# Not optional. You should see "SUCCESS! OpenCart successfully
# installed" go past; if you never see it, nothing below will work.
docker compose logs -f shop

# --- 3. Verify from outside, not from the logs ---------------------
# Three URLs must answer 200. waitfor.py checks all three and is the
# answer to "is my shop up?" for the rest of the semester.
python waitfor.py

# --- 4. Look at the shop, as a customer and as the owner -----------
# Storefront   http://localhost:8090/
# Admin        http://localhost:8090/admin/   (admin / admin123)
# Database     http://localhost:8091/         (shop / shoppw / opencart)

# --- 5. Count what you were given ----------------------------------
# In Adminer, or here. 159 tables; 19 products; 0 orders.
docker compose exec db mariadb -ushop -pshoppw opencart \
  -e "SELECT COUNT(*) FROM information_schema.tables \
      WHERE table_schema='opencart';"

# --- 6. Prove the volume is doing what you think -------------------
# Change something in the admin panel first -- edit a product name.
# Then stop WITHOUT -v and start again. Your change is still there.
docker compose down
docker compose up -d
python waitfor.py

# --- 7. Now reset it deliberately ----------------------------------
# -v throws the volumes away. The shop reinstalls from nothing and
# your edit is gone. Do this once, now, while it costs you nothing.
docker compose down -v
docker compose up -d
python waitfor.py

# --- 8. Measure it -------------------------------------------------
# Produces the table quoted above, on YOUR machine.
python m05_coldstart.py --warm
\end{lstlisting}

\smallskip
\textbf{What to notice.} Step~6 and step~7 differ by two characters and by
everything else. Having done both deliberately, in a week when the shop
contains nothing you care about, is the cheapest possible way to learn the
difference.

\smallskip
Also notice that step~3 does not trust step~2. The log said the install
succeeded; step~3 asks the shop. Throughout this course, the thing that
reports success and the thing that is true are kept apart on purpose.
\end{lab}

\begin{checkpoint}
\begin{tightnum}
  \item A classmate's shop returns ``500 Internal Server Error'' on every page.
        Which of the four layers in \dsref{lampstack} are known to be working,
        which is known to have failed, and where would you look first?
  \item Explain why \texttt{depends\_on} alone is not enough to order the
        database before the shop, and describe the failure it produces.
  \item You run \texttt{docker compose down -v} by mistake in Week~12. State
        exactly what is lost, what is not, and what would have had to exist
        beforehand for this to be a minor inconvenience.
\end{tightnum}
\end{checkpoint}

\begin{chaptersummary}
\begin{tight}
  \item A shop runs on four layers --- operating system, web server,
        application, database. Knowing which layer a symptom belongs to is most
        of diagnosis; only the database's loss is unrecoverable.
  \item An image is a built filesystem, a container is a running instance of
        one, and a volume is the only part that survives the container. A
        container is not a virtual machine, not a backup, and not secure by
        itself.
  \item The whole stack is three services in one Compose file: database, shop,
        database console.
  \item \texttt{depends\_on} orders starts, not readiness. A healthcheck is what
        makes the shop wait until the database will actually accept a login.
  \item \texttt{down} keeps your shop; \texttt{down -v} destroys it. Learn the
        difference in Week~3.
  \item The image is built from a pinned release verified by checksum, because
        no maintained OpenCart image exists to pull and because a version
        number alone is not reproducible.
  \item Measured on this machine: 338.4 s from scratch, 34.1 s cold, 16.4 s
        warm --- a ratio of 20.6$\times$ between the cost paid once and the cost
        paid daily. The stack occupies about 1.7 GB.
  \item The installed shop has 159 tables and 19 products, and eighteen of
        those tables are about products. One product is not one row.
\end{tight}
\end{chaptersummary}

\begin{reviewq}
  \item \textbf{(MCQ)} In the port mapping \texttt{8090:80}, the 80 is:
        (a)~the port on your machine; (b)~the port inside the container;
        (c)~the database port; (d)~an arbitrary label.
  \item \textbf{(MCQ)} Which survives \texttt{docker compose down}?
        (a)~Nothing; (b)~the containers; (c)~data written to a named volume;
        (d)~files edited inside a running container.
  \item \textbf{(Short)} Distinguish a container from a virtual machine, and
        give one consequence of the difference that matters in this course.
  \item \textbf{(Short)} Why does this course build the shop image from a pinned
        release rather than pulling a ready-made one? Give both reasons.
  \item \textbf{(Short)} A page loads but shows yesterday's price. Which layer
        is at fault, and what are the two most likely explanations?
  \item \textbf{(Applied)} Your stack will not start: port 8090 is already in
        use. Describe exactly what you would change, in which file, and what you
        would then have to change in the shop's own configuration for links to
        keep working. (Look at \texttt{OC\_HTTP\_SERVER} before answering.)
  \item \textbf{(Applied)} Design a measurement that would show whether the
        \emph{database} or the \emph{application} layer is responsible for a
        slow category page, using only tools in this chapter. State what you
        would record, at what catalogue sizes, and what result would point at
        each layer. \chref{catalogue} runs your experiment; write it first.
\end{reviewq}
